%\documentclass[a4paper,10pt]{scrreprt} %Dokumentklasse
%***********************************************************************************************%
%===============================================================================================%
% 																								%
%									    Bachelorarbeit zum Thema  								%
% 																								%
%  									   "Beispiel Abschlussarbeit" 								%
% 																								%
%									    von Gregor Zwickl, 2023									%
% 																								%
%===============================================================================================%
%***********************************************************************************************%


% ============= Klassendefinition ============= %

\documentclass[
	%draft, 		% Entwurfsstadium
	final, 			% fertiges Dokument
	paper=a4, 		% Papier Einstellungen
	pagesize=auto,	% Ist bereits default
	fontsize=12pt, 	% Schriftgröße
	ngerman, 		% Sprache
	openright,		% Neue chapter starten auf ungeraden(odd) Seitenzahlen, Alternative: open=right
	%twoside,		% Aktivieren, falls doppelseitiger Druck
	numbers=noendperiod,	% kein Punkt nach Nummerierung, eventuell autoendperiod
]{scrreprt}

% ============ Dokumentinformationen ============ %

\usepackage[
	pdftitle={Name Abschlussarbeit},		% Titel der Arbeit
	pdfsubject={Bachelorarbeit},			% Art der Abschlussarbeit, um was geht es
	pdfauthor={Tjalf Stadel},				% Name des Autors
	pdfkeywords={},							% Stichwörter
	pdfborder={0 0 0},						% {horizontalCornerRadius verticalCornerRadius borderWidth}
	colorlinks=true,						% Links sind farbig
	%breaklinks=true,						% unnötig
	linktocpage=true,						% Auch Seitenzahlen im Inhaltsverzeichnis können angeklickt 												% werden.
	citecolor=black,						% Farbe der Zitate
	menucolor=black,						% Farbe des ...
	urlcolor=black,							% Farbe der Links
	linkcolor=black,						% Farbe der Links
]{hyperref}							% mögl. Farben: (red, green, blue, cyan, magenta, yellow, black, white)


\title{Bachelorarbeit}
\author{Tjalf Stadel}
\date{02. Januar 2023}

% ============= Verwendete Pakete ============= %

\usepackage[T1]{fontenc} 				% Ausgabefont; bindet europäische (Sonder-)Zeichen ein
\usepackage[utf8]{inputenc} 			% Einbindung von Umlauten, Sonderzeichen (Standart ist nur ASCII)
% oder: \usepackage[latin1]{inputenc}
\usepackage{graphicx} 					% Einbinden von Bildern (.jpg; .png; .pdf)
%\usepackage[pdftex]{graphicx,color}
\usepackage[ngerman, english]{babel} 			% Deutsche Normen / Deutscher Sprachraum
\usepackage{geometry} 					% Seitenlayout
\usepackage{amsmath} 					% Matheumgebung  / Mathepacket
\usepackage{acronym}					% Abkürzungen erstellen
\usepackage[labelfont=bf]{caption}		% Beschriftungen (Schriftart=bf)
\usepackage{tocloft}					% Kontrolle über Layout (Inhaltsverzeichnis, Bilder, ...)
\usepackage{enumitem}					% Kontrolle über Layout (itemize, enumerate, description)
\usepackage{multirow}					% Zusammenfassen von Zeilen und Spalten in Tabellen
\usepackage{float}						% Defining floating objects like figures and tabels
%%\usepackage[headsepline,automark,nouppercase]{scrpage2}	%stellt vorgefertigte PageStyles zur Verfügung
\usepackage{scrlayer-scrpage}			% Kopf- und Fußzeilen mit KOMA-Script
\usepackage{pdfpages}					% PDF-Dokumente einbinden
\usepackage{tikz}						% Erstellen graphischer Elemente
\usepackage{pgfplots}					% Plots erstellen (normale oder logarithmische Skalierung)
\usepackage{pgfgantt}					% Gantt charts zeichen
\usepackage{pdflscape}					% landscapeorientierung von PDFs
\pgfplotsset{compat=newest}			% Kompatibilität mit neuerster Version (gibt immer wieder Updates)
\usepackage{grffile}					% Erweiterter file name support für Grafiken
\usepackage{subcaption}				% Typeset für sub-captions
\usepackage[nottoc]{tocbibind}			% fügt die Bib automatisch zum Inhaltsverzeichnis hinzu 													% (fügt Inhaltsverzeichnis nicht hinzu)
%\raggedbottom 							% bereits voreingestellt


%\usepackage{paralist} % Aufzälungsarten lassen sich ändern

%\usepackage{fancyhdr} % Kopf- und Fußzeile editierbar (nicht mit KOMA-Skript empfohlen)

%\pagestyle{fancy}	% verwendet package fancyhdr (nicht mit KOMA-Skript empfohlen)


% Einstellungen der Kopf- und Fußzeilen
\ohead{\headmark}									% Oben Rechts etwas angeben
\chead{}											% Keine Ausgabe oben in der Mitte
\cfoot[]{}											% Keine Seitenzahl in der Mitte
\ofoot[\pagemark]{\pagemark}						% Seitenzahlen rechts
\automark{section}									% Sektionen in Kopfzeile (Einseitiger Modus)
\renewcommand*{\headfont}{\normalfont}				% Kein kursiver Text in Kopfzeile
\pagestyle{scrheadings}

% Absätze und Abstände
\renewcommand{\baselinestretch}{1}\normalsize		% Zeilenabstand
\setlength{\parindent}{0pt}							% Kein Einrücken nach Absatz
\geometry{inner=31.5mm,outer=31.5mm,bottom=3cm}		% Abstand zu unterem Seitenrand
\textwidth147mm										% Breite des Textbereichs

%%%%%%%%%%%%%%%%%%%%%%%%%%%%%%%%%%%%%%%%%%%%%%%%%%%%%%%%%%%%%%
% ============= 	 Beginn des Dokuments  	   ============= %
%%%%%%%%%%%%%%%%%%%%%%%%%%%%%%%%%%%%%%%%%%%%%%%%%%%%%%%%%%%%%%




\begin{document}
\pagenumbering{Roman}							% Römische Seitenzahlen bis zum Hauptteil
% ============= Titelblatt ============= %
\pdfbookmark[0]{Titelseite}{title}				% Titelseite im pdf als Bookmark anzeigen
\includepdf[pages={1}]{images/deckblatt.pdf}
\addtocounter{page}{-1}

% ============= Leerseite ============ %
\newpage
\pagestyle{empty}
\phantom{t}										% Schafft Platz für Argument, zeigt dieses aber nicht an
\addtocounter{page}{-1}

% ============= Titelseite =========== %
\newpage
\begin{center}
	\begin{Large}
	\textbf{Titel der Arbeit auf DEUTSCH}\\
	\vspace{1cm}
	\textbf{Titel der Arbeit auf ENGLISCH}
\end{Large}

\vspace{4cm}

Bachelorarbeit\\
von\\
Muster Student

\vspace{4cm}

Durchgef\"uhrt am\\
\vspace{0.5cm}
\textbf{INSTITUT F\"UR FLUGZEUGBAU}\\
\textbf{UNIVERSIT\"AT STUTTGART}\\

\vspace{2cm}

Betreuer: M. Sc. Ruben Czichos

\vspace{2cm}

Stuttgart, im November 2018
\end{center}
\addtocounter{page}{-1}

% ==== Originale Aufgabenstellung ==== %
\includepdf{./images/aufgabenstellung.pdf}
\addtocounter{page}{-1}


% ============= Vorspann ============= %

\chapter*{Kurzfassung}
\addcontentsline{toc}{chapter}{Kurzfassung}	% Fügt Eintrag ins Inhaltsverzeichnis hinzu 
Lorem ipsum dolor sit amet, consetetur sadipscing elitr, sed diam nonumy eirmod tempor invidunt ut labore et dolore magna aliquyam erat, sed diam voluptua. At vero eos et accusam et justo duo dolores et ea rebum. Stet clita kasd gubergren, no sea takimata sanctus est Lorem ipsum dolor sit amet. Lorem ipsum dolor sit amet, consetetur sadipscing elitr, sed diam nonumy eirmod tempor invidunt ut labore et dolore magna aliquyam erat, sed diam voluptua. At vero eos et accusam et justo duo dolores et ea rebum. Stet clita kasd gubergren, no sea takimata sanctus est Lorem ipsum dolor sit amet. Lorem ipsum dolor sit amet, consetetur sadipscing elitr, sed diam nonumy eirmod tempor invidunt ut labore et dolore magna aliquyam erat, sed diam voluptua. At vero eos et accusam et justo duo dolores et ea rebum. Stet clita kasd gubergren, no sea takimata sanctus est Lorem ipsum dolor sit amet.\\   
Duis autem vel eum iriure dolor in hendrerit in vulputate velit esse molestie consequat, vel illum dolore eu feugiat nulla facilisis at vero eros et accumsan et iusto odio dignissim qui blandit praesent luptatum zzril delenit augue duis dolore te feugait nulla facilisi. Lorem ipsum dolor sit amet, consectetuer adipiscing elit, sed diam nonummy nibh euismod tincidunt ut laoreet dolore magna aliquam erat volutpat.\\ 
Ut wisi enim ad minim veniam, quis nostrud exerci tation ullamcorper suscipit lobortis nisl ut aliquip ex ea commodo consequat. Duis autem vel eum iriure dolor in hendrerit in vulputate velit esse molestie consequat, vel illum dolore eu feugiat nulla facilisis at vero eros et accumsan et iusto odio dignissim qui blandit praesent luptatum zzril delenit augue duis dolore te feugait nulla facilisi.\\   
Nam liber tempor cum soluta nobis eleifend option congue nihil imperdiet doming id quod mazim placerat facer possim assum. Lorem ipsum dolor sit amet, consectetuer adipiscing elit, sed diam nonummy nibh euismod tincidunt ut laoreet dolore magna aliquam erat volutpat. Ut wisi enim ad minim veniam, quis nostrud exerci tation ullamcorper suscipit lobortis nisl ut aliquip ex ea commodo consequat.

\newpage
\chapter*{Abstract}
\addcontentsline{toc}{chapter}{Abstract}
\newpage
%\pagestyle{scrheadings}
\renewcommand*{\chapterheadstartvskip}{\vspace*{2.3\baselineskip}}	% Abstand einstellen
%\pdfbookmark[0]{Inhaltsverzeichnis}{toc}							% Inhaltsverzeichnis im pdf als 																			% Bookmark anzeigen
\renewcommand{\baselinestretch}{1.2}\normalsize						% Zeilenabstand
\tableofcontents													% Inhaltsverzeichnis
\renewcommand{\baselinestretch}{1.25}\normalsize					% Zeilenabstand
\renewcommand*{\chapterheadstartvskip}{\vspace*{2.3\baselineskip}}	% Abstand einstellen
% ============= Symbolverzeichnis ============= %
\chapter*{Nomenklatur}
\addcontentsline{toc}{chapter}{Nomenklatur}
Zur Förderung der Übersichtlichkeit werden in der vorliegenden Arbeit Skalare mit normaler Strichstärke und Vektoren bzw. Matrizen fettgedruckt dargestellt.
\renewcommand{\baselinestretch}{1.1}\normalsize						% Zeilenabstand
\begin{table}[H]
	\begin{tabular}{ll}
		\hline
		\textbf{Symbol}\phantom{123456789}& \textbf{Bedeutung} \phantom{123456789123456789123456789123456789}\\ 
		\hline
		\\
		Kräfte\\
		\hline
		$A$					& Auftrieb\\
		$F$					& Schub\\
		$W$					& Widerstand\\
		$Y$					& Seitenkraft\\
		$R$					& Reibung\\
		$F_{\ddot{U}berschuss}$ & Überschussleistung\\
		$W_{min}$			& minimaler Widerstand\\
		\\
		Momente\\
		\hline
		$L$					& Rollmoment\\
		$M$					& Nickmoment\\
		$N$					& Giermoment\\
		\\
		Aerodynamische Beiwerte\\
		\hline
		$C_{A}$				& Auftriebsbeiwert\\
		$C_{W}$				& Widerstandsbeiwert\\
		$C_{M}$				& Nickmomentenbeiwert\\
		$C_{A,max}$			& maximaler Auftriebsbeiwert\\
		$C_{A,roll}$		& Auftriebsbeiwert beim Rollvorgang\\
		$C_{AP,0}$			& Profilauftriebsbeiwert bei $\alpha = 0$\\
		$C_{A0}$			& Auftriebsbeiwert bei $\alpha = 0$\\
		$C^*_A$				& Auftriebsbeiwert im Punkt des besten Gleitens\\
		$C_{W,roll}$		& Widerstandsbeiwert beim Rollvorgang\\
		$C_{W,profil}$		& Widerstandsbeiwert des Profils\\
		$C_{W0}$			& Nullwiderstandsbeiwert (Widerstandsbeiwert bei $C_A=0$)\\
		$C_{Wi}$			& induzierter Widerstandsbeiwert\\
		$C_{W,min}$			& minimaler Widerstandsbeiwert\\
		$C_{WP,A}$			& auftriebsabhängiger Profilwiderstandsbeiwert\\
		$C^*_W$				& Widerstandsbeiwert im Punkt des besten Gleitens\\
%		$C_{L}$				& Rollmomentenbeiwert\\
%		$C_{N}$				& Giermomentenbeiwert\\
%		$C_{Y}$				& Seitenkraftbeiwert\\
		
		$C_{a}$				& Auftriebsbeiwert des Profils\\
		$C_{w}$				& Widerstandsbeiwert des Profils\\
		$C_{m}$				& Nickmomentenbeiwert des Profils\\
%		$C_{l}$				& Rollmomentenbeiwert des Profils\\
%		$C_{n}$				& Giermomentenbeiwert des Profils\\
%		$C_{y}$				& Seitenkraftbeiwert des Profils\\
		\\
		Geschwindigkeiten\\
		\hline
		$V$					& Geschwindigkeit\\
		$V_{min}$			& Minimalgeschwindigkeit\\
		$V_{Start}$			& Abhebegeschwindigkeit\\
		$V^*$				& Geschwindigkeit im Punkt des besten Gleitens\\
		$u$					& horizontale Geschwindigkeit (in x-Richtung)\\
		$V_h$				& horizontale Gewchwindigkeitskomponente\\
%		$v$					& Geschwindigkeit in y-Richtung\\
		$w$					& Sinkgeschwindigkeit\\
		$V_v$				& Steiggeschwindigkeit / vertikale Geschwindigkeitskomponente / SEP\\
		$V_K$				& Kurvengeschwindigkeit\\
		$V_{K,min}$			& minimaler Kurvengeschwindigkeit\\
		$a$					& Schallgeschwindigkeit\\
		$Ma$				& Machzahl\\
		$p$					& Drehgeschwindigkeit um die Rollachse\\
		$q$					& Drehgeschwindigkeit um die Nickachse\\
		$r$					& Drehgeschwindigkeit um die Gierachse\\
		\\
	\end{tabular}
\end{table}
\begin{table}[H]
	\begin{tabular}{ll}
		Geometrische Größen\\
		\hline
		$S$					& Bezugs(flügel)fläche\\
		$l_\mu$				& Bezugsflügeltiefe\\
		$l$					& Profiltiefe, Länge\\
		$b$					& Spannweite\\
		$\Lambda$			& Flügelstreckung\\
		$\lambda$			& Zuspitzung\\
%		$s$					& Halbspannweite\\
		$r_K$				& Kurvenradius\\
		$r_{K,min}$			& minimaler Kurvenradius\\
		$s$					& Flugstrecke\\
		$h$					& Höhe\\
		\\
		Aerodynamische Größen\\
		\hline
		$Re$				& Reynoldszahl\\
		$\alpha$			& Anstellwinkel\\
		$\epsilon$			& Gleitzahl\\
		$\gamma$			& Bahnneigungswinkel\\
		$\chi$				& Azimutwinkel\\
		$\phi$				& Rollwinkel\\
		$\rho$				& Luftdichte\\
		$T$					& Temperatur\\
		$\eta$				& dynamische Viskosität\\
		$\nu$				& kinematische Viskosität\\
		$p$					& Druck\\
		$q$					& Staudruck\\
		$k$					& Widerstandsfaktor\\
		$k_i$				& Widerstandsfaktor induzierter Widerstand\\
		$k_P$				& Widerstandsfaktor auftriebsabhängiger Profilwiderstand\\
		$\Phi_A$			& Einflussfaktor des Bodeneffekts auf den Auftrieb\\
		$\Phi_W$			& Einflussfaktor des Bodeneffekts auf den Widerstand\\
		$\beta_A$			& Faktor für $\Phi_A$\\
		$\beta_W$			& Faktor für $\Phi_W$\\
		$\delta_A$			& Faktor für $\Phi_A$\\
		$\delta_W$			& Faktor für $\Phi_W$\\
		\\
		Weitere Größen\\
		\hline
		$m$					& Masse\\
		$g$					& Gewichtskonstante\\
		$n$					& Lastvielfaches\\
		$n_{max}$			& maximales Lastvielvaches\\
		$t$					& Zeit\\
		$\mu_{roll}$		& Rollreibungskoeffizient\\
%		$\dot{m}$			& Massendurchsatz\\
	\end{tabular}
\end{table}




\begin{table}[H]
	\begin{tabular}{ll}
		\hline
		Symbol\phantom{123456789}& Bedeutung \phantom{123456789123456789123456789123456789}\\ 
		\hline
		$\boldsymbol{A}$	& Prozessmatrix eines Systems / Beschleunigungsvektor\\
		$\boldsymbol{B}$	& Eingangsmatrix eines Systems\\
		$\boldsymbol{C}$	& Beobachtungsmatrix eines Systems\\
		$\boldsymbol{H}$	& Approximierte Messmatrix\\
		$\boldsymbol{P}$	& Schätzfehlerkovarianzmatrix\\
		$\boldsymbol{Q}$	& Prozessrauschkovarianzmatrix\\
		$\boldsymbol{R}$	& Messrauschkovarianzmatrix / Sichtlinie\\
		$\dot{\boldsymbol{R}}$	& Relativgeschwindigkeit\\
		$\ddot{\boldsymbol{R}}$	& Relativbeschleunigung\\
		$\boldsymbol{S}$	& Innovationskovarianzmatrix\\
		$\boldsymbol{T}_{\alpha}$	& Drehmatrix\\
		$\boldsymbol{Z}$	& Filterinnovation\\
		$\boldsymbol{\phi}$	& Jacobi-Matrix\\
		$\boldsymbol{\Pi}$	& Übergangsmatrix\\
		$M$					& Anzahl Monte-Carlo-Simulationen\\
		$N$					& Navigationskonstante / Filteranzahl\\
	\end{tabular}
\end{table}

\begin{table}[H]
	\begin{tabular}{ll}
		\hline
		Abkürzung\phantom{123456}& Bedeutung \phantom{1234567891234567891234567891234567899}\\ 
		\hline
	\end{tabular}
\end{table}
\begin{addmargin}[0.2cm]{0cm}
\vspace{-0.6cm}
\begin{acronym}[MMAE~~~~~~~~~~~~]
	\setlength{\itemsep}{-\parsep} 
	\acro{APN}{Augmented Proportional Navigation}
	\acro{CA}{Constant Acceleration}
	\acro{CV}{Constant Velocity}
	\acro{CM}{Cruise Missiles}
	\acro{CT}{Constant Turn}
\end{acronym}
\end{addmargin}
\renewcommand{\baselinestretch}{1.25}\normalsize					% Zeilenabstand
%\markleft{Nomenklatur}

\renewcommand*{\listoffigures}{
	\begingroup
	%\tocchapter
	\tocfile{\listfigurename}{lof}
	\endgroup
}

\cleardoublepage
\setlength{\cftfigindent}{0cm}
\renewcommand{\baselinestretch}{1.1}\normalsize						% Zeilenabstand
%\pdfbookmark{Abbildungsverzeichnis}{lof} 
\listoffigures 														% Abbildungsverzeichnis
%\markleft{Abbildungsverzeichnis}

\renewcommand*{\listoftables}{
	\begingroup
	%\tocchapter
	\tocfile{\listtablename}{lot}
	\endgroup
}
%\pdfbookmark[0]{Tabellenverzeichnis}{tbv}							% Tabellenverzeichnis im pdf als Bookmark anzeigen
\setlength{\cfttabindent}{0em}
\listoftables														% Tabellenverzeichnis
\renewcommand{\baselinestretch}{1}\normalsize						% Zeilenabstand

% ============= Hauptteil ============= %

\pagenumbering{arabic}  											% Arabische Zahlen im Hauptteil bis zum Anhang

\selectlanguage{ngerman}                % Set language to german
%********************************
\chapter{Einleitung}					% Kapitel
\label{chap:einleitung}					% Querverweis auf "einleitung"kann aufgerufen werden
%********************************		% mit: \ref{chap:einleitung}

%********************************
\section{Motivation}					% Sektion
\label{sec:motivation}
%********************************
Hier erfolgt die Hinführung zum Thema. Worum geht es? Was ist der Sinn und Zweck dieser Arbeit? Dem Leser wird hier auch erklärt, was ihn wo in dieser Arbeit erwartet.

Diese Bachelorarbeit ist entsprungen, um in der Akamodell Stuttgart effizienter Modellflugzeuge auslegen zu können. Die Akamodell Stuttgart ist ein eingetragener Verein an der Universität Stuttgart, an dem Studenten lernen Modellflugzeuge zu entwerfen und zu bauen. Dazu nimmt die Akamodell zwei-jährlich an der Air Cargo Challenge (ACC) teil. Die ACC ist ein Wettbewerb, bei dem verschiedene Teams aus europäischen und einigen internationalen Universitäten gegeneinander antreten, um ein Flugzeug zu konstruieren, dass dem Regelwerk entsprechend die Flugaufgaben möglichst gut zu erfüllen. Besonders für diese Wettbewerbe ist eine schnelle und qualitativ hochwertige Auslegung von Bedeutung. Da solche Auslegungen viel Know-how erfordern, ist es das Ziel ein Tool zu entwickeln, dass dieses vereinfacht.

\subsection{Motivation genau}			% Untersektion
\label{subsec:motgenau}

Im Speziellen befasst sich diese Arbeit mit der Massenabschätzung des Flugzeuges.

%********************************
\section{Zielsetzung}
\label{sec:zielsetzung}
%********************************
In dieser Arbeit soll ein Algorithmus entwickelt werden, der es ermöglicht  in jedem Schritt der Modellflugzeugauslegung 

%********************************
\section{Gliederung der Arbeit}
\label{sec:gliederung}
%********************************
In Folgenden werden zunächst die Grundlagen und Grundlagen erklärt. Danach wird der Stand der Technik der Flugzeugauslegung erklärt und der Bau und Aufbau von Modellflugzeugen erläutert. Danach 


\selectlanguage{ngerman}
%********************************
\chapter{Grundlagen}
\label{chap2:grundlagen}
%********************************
%********************************
\section{Stand der Technik}
\label{sec:sdt}
%********************************
Was sind die neusten Erkenntnisse auf dem Gebiet?

%********************************
% Schreiben von unterschieden, Motorisiert/Segelflug
\section{Unterschiede zwischen Seglern und motorisierten Flugzeugen}

Einen Teil im Modellflug macht das Segelfliegen aus. Diese haben einige Vorteile gegenüber motorisierten Fliegern. Jene sparen Gewicht, da kein Motor und nur eine kleinere Batterie benötigt wird. Dafür können diese nur 


\section[short]{Aufgaben}

\subsection*{Startstrecke}
Dieses Trifft nur zu, wenn das Flugzeug aus eigener Kraft starten soll und nicht geworfen oder geschleppt wird.

Je nach Gelände auf dem Geflogen werden soll, kann die Länge der Startstrecke von Relevanz sein. Probleme stellen dabei die Unterschiedlichen Rollwiederständen verschiedener Untergründe. Dabei hat Asphalt einen deutlich niedrigeren als hohes Gras. Für einen schnellen Start ist daher 

\section{Air Cargo Challenge}

Im Folgenden werden einige Grundlagen erklärt. Dieses wird anhand des Beispiels ACC2022 getan.

\subsection{Konstruktionsbeschränkungen}

Im Modellflug ist die Größe des Modellflugzeuges hauptsächlich durch die EU-Drohnenverordnung geregelt. Dadurch dürfen nur Drohnen unter 25 kg geflogen werden, solange kein EU-Fernpiloten-Zeugnis vorhanden ist. Zusätzlich werden im ACC Regelwerk noch weitere Beschränkungen geregelt. Der Motor sowie die Batterie sind vorgegeben und das Flugzeug muss in eine Box mit den Seitenlängen $1.5\;\text(m)$ passen und eine maximal höhe von $1$ nicht überschreiten


\subsection{Flugaufgabe}

Für jedes Modellflugzeug muss eine Flugaufgabe definiert werden, auf die optimiert werden kann. Für die ACC muss zunächst eine möglichst kurze Startstrecke von 100 Meter bzw 50 Meter erreicht werden.

%********************************
\chapter{Hauptteil}

text

%********************************
\chapter{Ergebnisse}
\label{chap4:ergebnisse}
%********************************
Hier werden alle während der Arbeit gewonnenen Informationen ausführlich dargestellt, diskutiert und interpretiert. Dabei sollte im Bezug zur Ausgangsituation ein Ergebnis/Erkenntnisgewinn bzw. eine Verbesserung/Änderung vorliegen und eine Schlussfolgerung mit Empfehlung gezogen werden. Aufgetretene Diskrepanzen und mögliche Ursachen sollten klar dargestellt werden.

%********************************
\chapter{Zusammmenfassung und Ausblick}
\label{chap:zusammenfassung_und_ausblick}
%********************************
Zum Schluss sollte nochmals die wesentlichen Ergebnisse klar herausgestellt werden. Außerdem kann ein Ausblick für weitere Arbeiten/Untersuchungen in diesem Bereich erfolgen. Umfang beträgt hier ca. 2 Seiten.


% ============= Literaturverzeichnis ============= %
%\renewcommand{\baselinestretch}{1.2}\normalsize						% Zeilenabstand
\clearpage
\bibliographystyle{plaindin}											% Art von Seitenumbruch
\addcontentsline{toc}{chapter}{Literaturverzeichnis}
\bibliography{Literatur}
%\renewcommand{\baselinestretch}{1}\normalsize						% Zeilenabstand
% =============Anhang ============= %
\appendix
\selectlanguage{ngerman}
%********************************
\addchap{Anhang}
\refstepcounter{chapter} 
\label{chap:anhang}
Hier können noch Tabellen, Messprotokolle, Rechnerprotokolle, Konstruktionszeichnungen sowie Programmcodes und ähnliches zur Dokumentation angeheftet werden.

% ============= Nutzungsrechte ========= %
\newpage
\pagestyle{empty}
\input{chapters/98_Nutzungsrechte}

% ============= Leerseite ============= %
\newpage
\pagestyle{empty}
\phantom{t}


%%%%%%%%%%%%%%%%%%%%%%%%%%%%%%%%%%%%%%%%%%%%%%%%%%%%%%%%%%%%%%
% ============= 	  Ende des Dokuments  	   ============= %
%%%%%%%%%%%%%%%%%%%%%%%%%%%%%%%%%%%%%%%%%%%%%%%%%%%%%%%%%%%%%%
\end{document}